\documentclass[11pt]{article}
\usepackage[utf8]{inputenc}
\usepackage[T1]{fontenc}
\usepackage{lmodern}
\usepackage[margin=1in]{geometry}
\usepackage{amsmath,amssymb}
\usepackage{textcomp}
\usepackage[dvipsnames]{xcolor}
\usepackage{enumitem}
\usepackage{parskip}
\usepackage{titlesec}
\usepackage[colorlinks=true,linkcolor=blue,citecolor=blue,urlcolor=blue]{hyperref}
\usepackage{natbib}

\titleformat{\section}{\large\bfseries}{\thesection}{1em}{}
\titleformat{\subsection}{\normalsize\bfseries}{\thesubsection}{1em}{}
\titleformat{\subsubsection}{\normalsize\bfseries\itshape}{}{0em}{}

\title{\textbf{The Channel Is Not the Content:\\ A Calibrated Decomposition of Capability, Policy, and Format Drift across LLM Version Pairs}}
\author{}
\date{}

\newcommand{\ph}{\textit{[placeholder --- fill in after experiments are run]}}

\begin{document}
\maketitle
\vspace{-2em}

\begin{quote}
\small\textit{Assembled from the six confirmed stage outputs of an \texttt{auto-research-discussion} run (seed topic B3.2), post research-taste audit (\textcolor{ForestGreen}{\textbf{PASS}}), post stage review (two REVISE fixes applied --- a nearest-neighbor sentence in Related Work, and per-RQ falsification conditions in Experiments). Revised 2026-08-21 per the 2026-08 novelty/feasibility review (F1 2026 repositioning, F2 mixed-axis/scale-ladder calibration bounds, F3 RQ6 rank-unit fix, F4 off-ceiling pre-verification). Status: \textbf{specification / paper draft --- no experiments have been run.} All numbers below are pre-registered design targets, not results.}
\end{quote}

\begin{abstract}
When an LLM provider ships a new model version, downstream users observe behavioral change but are rarely told what actually moved underneath: the model's capability, its alignment or safety policy, or merely its default output formatting. Standard longitudinal audits report a single blended before/after delta that cannot distinguish these three causes --- and the canonical audit of ChatGPT's behavior over time found that part of an apparent capability regression was in fact a formatting change (added markdown fencing) that broke an executability check, misattributing surface drift as capability loss. We decompose behavioral drift between LLM version pairs into capability, policy, and format components using three probe batteries, each engineered to hold the other two axes fixed by construction, an explicit cross-axis leakage estimate, and a within-version noise floor that separates ordinary sampling variability from genuine drift. Because no real version pair carries a verified ground-truth cause, we calibrate the entire procedure by testing blind recovery of the known dominant axis on constructed control pairs before interpreting any real pair's decomposition. Applied to version pairs from multiple vendors, we test whether the resulting axis-dominance structure recurs across providers or is vendor-specific, and whether format drift --- invisible to blended audits --- carries a non-negligible share of what current practice reports as capability change.
\end{abstract}

\section{Introduction}

Large language model providers now ship new model versions frequently, often within weeks of one another, and the practice of building on top of a pinned API endpoint or a downloaded checkpoint has become routine enough that a version bump is treated as background noise rather than an event worth scrutinizing on its own terms. Users and integrators do notice behavioral change after such an update --- a task that used to succeed now fails, a response that used to comply now refuses, an output that used to parse cleanly now doesn't --- but the standard way of reporting this change, a single blended before/after score, treats all three of these very different events identically.

Yet a central question about every such update goes unanswered by the number itself: did the model actually get smarter, did its alignment or safety policy shift, or did nothing change except its default formatting habits? These three have completely different, and completely non-interchangeable, downstream consequences. A capability change calls for re-benchmarking task performance; a policy change calls for re-auditing safety and compliance posture; a format change calls for nothing more than fixing a brittle parser. Conflating them is not a hypothetical risk. The canonical longitudinal audit of ChatGPT's behavior over time \citep{ChatGPTBehavior2307} found that part of an apparent capability regression in code generation was in fact caused by newly added markdown fencing breaking an executability check --- a pure formatting change that got scored and reported as a capability loss. Nor is this aggregate-number-hides-separable-causes structure specific to version auditing: a fine-grained analysis of how LLMs generalize under SFT versus RL post-training shows that a single blended accuracy delta conceals a decomposable shift in the underlying composition of reasoning skills \citep{HowandWhyLLMsGeneralizeA4063} --- the same conflation, arising there along a training-method axis rather than a version axis. If today's standard audits cannot reliably rule this out, every headline claim that a model's capability degraded between snapshots is potentially confounded by exactly the surface artifact this paper is built to separate out.

The direct predecessor of this problem is the longitudinal audit that first documented it: the ChatGPT behavior-over-time study \citep{ChatGPTBehavior2307} reports a single blended per-task delta between snapshots, with no mechanism to attribute that delta to capability, policy, or format, and its own reported code-generation drop is direct evidence of the conflation this paper targets --- a result worth stating plainly before surveying the wider landscape, since the paper whose central limitation motivates this one is also where the phenomenon was first put on record. The nearest \emph{contemporary} neighbors, however, are a fast-moving 2026 cluster on model-update governance and change detection that now surrounds this problem without occupying it: deployer-side update governance tests a new version by risk category --- functionality, formatting, and safety among them --- and explicitly names detecting and explaining model drift under limited provider transparency as an open challenge \citep{TestBeforeYouDeploy7789}; black-box change detection establishes \emph{that} a served model changed, token-efficiently, but not why \citep{TokenEfficientChangeDet1083}; item-level reliable-change analysis detects \emph{which} items genuinely moved against a within-model sampling-variability floor --- anticipating the noise floor used here, a debt we acknowledge explicitly --- but stops at binary change rather than cause \citep{BeyondTheMean7405}; and anytime-valid drift attribution in evaluation pipelines apportions a metric shift between the serving system and the judge, a different pair of causes anchored by human labels \citep{WhoDrifted5474}. Detection, per-category gating, and two-way system-vs-judge attribution are thus all occupied; a black-box, calibration-gated three-way attribution of a version pair's drift to capability, policy, or format is not. Beyond that cluster, related work approaches the problem from adjacent but distinct angles. Work on stage-wise alignment impact \citep{MoreRLHFMoreTrustFpiC} attributes behavioral change to specific training stages using influence-function-based data attribution, but this requires white-box access to checkpoints and training data that commercial version pairs never expose, and it measures only trustworthiness dimensions with no capability or format axis. Instruments built on per-token divergence \citep{FewTokensDeep6Mxh} or on weight differences \citep{DiffInterpretation6As4} between two models can isolate a single axis with fine granularity, but both require logit or full-weight access unavailable for API-served commercial pairs, and neither apportions a total observed difference across competing causes. A rank-correlation measure separating refusal propensity from knowledge limitation \citep{RefusalIndex9gJB} demonstrates that policy and capability are distinct, separately measurable constructs --- but only within a single model at one point in time, with no version-pair or drift framing at all. And judge-based evaluation frameworks \citep{WildBenchMKEH} document that surface length and style systematically distort automated scoring, treating this as a nuisance bias to be penalized out of a ranking metric rather than as a first-class, separately reported axis of behavioral change.

Turning this gap into a measurable question surfaces four separate obstacles. There is no ground truth for which cause produced a given real version pair's behavioral difference, so any attribution claim risks being unfalsifiable unless it can be checked indirectly, on pairs where the cause happens to be known. Probes meant to isolate one axis can still be contaminated by another in ways invisible from the probe's own output --- a policy change can lengthen or shorten a response in a way that looks like format drift, and a capability change can flip a refusal decision in a way that looks like policy drift. Ordinary run-to-run sampling variability, and the possibility of silent infrastructure changes mid-study, can masquerade as genuine drift unless it is measured and subtracted rather than assumed away. And because vendors may implement conceptually similar changes --- a tightened refusal rule, say --- through entirely different underlying mechanisms, from a policy layer to a system-prompt edit to an output filter, any claim that an attribution structure is shared across vendors risks reflecting how the instrument happened to be built rather than a real fact about the industry.

Our approach turns each of these obstacles into a specific, checkable design choice rather than a hedge. Because output text is simultaneously the only channel available for black-box measurement and itself a target of vendor change, while content changes do not alter surface-feature scores the way format changes do, the three probe batteries are each engineered to hold the other two axes fixed by construction --- freezing the channel to recover the content. Because no real pair carries a known cause, the entire pipeline is calibrated blind on constructed control pairs whose dominant axis is known by design, and no claim about a real pair is made unless that calibration gate passes first. Cross-axis contamination is not assumed away but explicitly estimated and reported as its own leakage share, and a within-version noise floor --- validated for stability across interleaved query windows --- determines how large an axis's contribution must be before it counts as drift at all rather than sampling noise. Finally, cross-vendor commensurability is tested empirically, by applying the identical instrument to every vendor and checking whether the resulting axis-dominance ordering agrees across them more than chance would predict.

We evaluate this design against version pairs from multiple commercial and open-weight vendors, alongside a blended-delta baseline that reproduces standard current practice and a same-model null-pair placebo that must show no detectable ``drift'' at all.

In summary, we contribute:
\begin{itemize}[leftmargin=1.4em]
\item A calibration protocol that validates a behavioral drift-attribution procedure on constructed control pairs of known cause before any claim is made about real, black-box version pairs where the true cause is unverifiable.
\item An explicit cross-axis leakage estimate and within-version noise-floor model that turn ``the probe batteries are orthogonal'' and ``this delta is really drift'' from silent assumptions into reported, falsifiable quantities.
\item A cross-vendor structural test of whether capability/policy/format axis-dominance is a shared industry pattern or a vendor-idiosyncratic one, evaluated behaviorally rather than assumed from vendor documentation that does not exist.
\item An empirical account of whether format drift --- invisible to a blended before/after score --- carries a non-negligible share of what standard longitudinal audits currently report as capability or policy change.
\end{itemize}

\section{Problem Definition}

\noindent\textbf{Formal setup.} Consider a version pair $P = (v_{\mathrm{old}}, v_{\mathrm{new}})$ from the same provider, both accessible only through ordinary text input/output --- no logits, weights, or training data at any point in the pipeline. A pair is defined as two directly comparable, adjacent shipped snapshots, simultaneously queryable over the study window, not a chain spanning multiple intermediate releases. Behavioral change is measured along three candidate sources $a \in \{\text{cap}, \text{pol}, \text{fmt}\}$, using probe batteries whose items are queried identically against both versions with $k = 5$ repeated samples per item per version at fixed decoding settings. For each source $a$, an axis share $\eta^2_a(P)$ is the proportion of the observed between-version difference attributable to that source, estimated in excess of a within-version noise floor $\sigma^2_w(P)$ computed from the repeated samples. A leakage share $\lambda(P)$ captures the proportion of one axis's measured signal explained by another axis's independently measured features; it is reported as its own quantity, never silently folded into $\{\eta^2_a\}$. The dominant axis of a pair is $\mathrm{dom}(P) = \arg\max_a \eta^2_a(P)$. For a set of pairs constructed so the true dominant axis is known in advance, recovery accuracy $\mathrm{Acc}$ is the fraction for which $\mathrm{dom}(P)$ matches the known cause. Cross-vendor agreement is summarized by Kendall's $W$ computed over per-pair rankings of the three axis shares --- each of $m \geq 6$ pairs contributes one ranking of the axes as a ranking unit --- tested against a permutation null; vendor-level structure is reported descriptively via clustering on the axis-share simplex, because with only three vendors as ranking units the exact permutation null cannot reach $p < 0.05$ except at perfect agreement.

\noindent\textbf{Why this is hard.} No real version pair carries a verified ground-truth cause of its behavioral difference, so an attribution claim about $\eta^2_a(P)$ is unfalsifiable unless the underlying procedure's reliability can be checked some other way --- specifically, on constructed pairs where the true cause is known, which is what licenses (or refutes) trusting the same procedure on real pairs where it is not. The three sources are not orthogonal in raw behavior: a policy change can alter response length in a way that looks like format drift, and a capability change can flip a refusal decision in a way that looks like a policy change, so any reported $\eta^2_a(P)$ risks silently absorbing another axis's effect unless the leakage between axes is estimated and reported rather than assumed to be zero. Because LLM outputs are stochastic and serving infrastructure can shift mid-study without notice, a genuine between-version difference must first be distinguished from ordinary within-version sampling variability, or $\sigma^2_w(P)$ itself becomes an unreliable reference and every downstream $\eta^2_a(P)$ estimate is corrupted along with it. Finally, because vendor-internal mechanisms are never directly observable, whether ``capability,'' ``policy,'' and ``format'' are commensurable constructs across providers cannot be verified from documentation or source access --- it can only be tested behaviorally, by checking whether an identical instrument applied to every vendor's pairs produces a rank ordering of axis dominance that agrees across vendors more than a permutation null would. Two further identifiability limits are intrinsic to the calibration strategy itself and are reported as first-class limitations rather than hidden: the constructed control pairs are defined under the same capability/policy/format taxonomy they are used to validate, so passing calibration establishes that taxonomy-aligned interventions are recoverable --- not that the taxonomy exhaustively spans real drift; and control pairs are necessarily built from open-weight models while the audited targets include commercial-scale serving stacks, so calibration transfer across scale and stack is an assumption to be stress-tested (via the mixed-axis controls and scale ladder below), not a given.

\section{Method Overview}

At a conceptual level, the study is a single measurement pipeline with four cooperating parts. Probe batteries and constructed control pairs enter the pipeline; a calibration step establishes whether the resulting decomposition can be trusted at all before any real pair is interpreted; a leakage-and-noise accounting model turns the raw per-item behavioral deltas into reportable axis shares with their own uncertainty and contamination bounds attached; and a cross-vendor agreement test turns the resulting per-pair decompositions into a claim about whether the same structure recurs across providers.

The \textbf{axis-isolating probe batteries and control-pair calibration} module builds three probe batteries, each engineered so the other two axes have minimal effect on its measured quantity by construction, and separately constructs control pairs with a known dominant axis, running the full pipeline blind on them to check whether it recovers the true cause. The \textbf{cross-axis leakage estimation} module regresses each axis's measured signal on the other axes' independently measured features, reporting the explained fraction as an explicit leakage share rather than letting it silently bias the reported axis shares. The \textbf{noise-floor variance model} estimates within-version run-to-run variance from repeated sampling and fits a mixed-effects model so an axis's contribution to observed drift counts only in excess of that floor, while also checking that the floor itself is stable across the study window. The \textbf{cross-vendor structural agreement test} summarizes each pair as a point on the axis-share simplex and tests whether the per-pair rank ordering of axis dominance agrees across pairs --- and, descriptively, across vendors --- more than a permutation null would predict.

\section{Methodology}

\subsection{Axis-Isolating Probe Batteries and Control-Pair Calibration}
\textbf{Challenge.} No real version pair carries a known, verifiable cause for its behavioral difference, so any attribution claim risks being unfalsifiable --- there is nothing to check it against directly.

\textbf{Insight.} Entanglement cannot be resolved analytically on real pairs, because no real pair offers ground truth --- but a decomposition procedure's \emph{reliability} can still be checked indirectly, by constructing pairs where the true cause is known by design and running the exact same batteries and model on them, blind, before ever applying the procedure to a real pair whose cause is unknown.

\textbf{Method.} Three probe batteries are each engineered so the other two axes have minimal effect on its measured quantity by construction: the capability battery uses format-robust extraction (executed code, normalized numeric matching, constrained multiple-choice extraction) so its score is invariant to surface formatting; the format battery is restricted to prompts on which both versions are already near-ceiling correct, so capability differences cannot leak into the format measurement; the policy battery stratifies prompts by independently known ground-truth harmfulness, so a refusal-rate change can be read against a fixed reference rather than confounded with correctness. Separately, control pairs are constructed by modifying exactly one axis at a time --- a base model against the same model fine-tuned on domain data for a capability control, a base model against a safety-tuned variant (following the safety-alignment fine-tuning recipe of \citealp{ContrastiveOverRefusalGXcN}) for a policy control, and identical weights under an altered system template for a format control --- so that, unlike any real vendor pair, the true dominant axis is known. The complete, unmodified pipeline is then run with pair identities withheld from the analysis code, and the predicted dominant axis is compared against the true one.

Two design elements bound what this calibration can and cannot license, and both are treated as part of the protocol rather than as afterthoughts. First, single-axis controls are constructed under the same three-axis taxonomy they are meant to validate, so blind recovery on them shows that the instrument detects taxonomy-aligned interventions --- not that real drift is spanned by the taxonomy --- and real vendor pairs are plausibly mixtures of causes, not pure cases. The control set therefore also includes mixed-axis pairs: pairs constructed by applying two interventions at a known relative dose (for example, a safety-tuning update combined with a system-template change), whose known dominant axis must be recovered under the same blind protocol, so that calibration covers the mixed regime in which the paper's headline real-pair claims actually live. Second, control pairs are built from small open-weight models while the audited pairs include commercial-scale serving stacks, so calibration transfer is itself an assumption rather than a consequence of the design: a scale ladder repeats the format-only control and the null-pair placebo at two or more open-weight scales and checks whether recovery accuracy and leakage share are stable in scale. A strong scale trend would flag that calibration numbers cannot be carried across to frontier pairs; whichever way it comes out, the transfer step is reported as an explicit limitation, never assumed silently.

\subsection{Cross-Axis Leakage Estimation}
\textbf{Challenge.} A probe battery nominally isolating one axis can still be contaminated by another axis in a way that is invisible from that battery's own output alone.

\textbf{Insight.} Because an axis that leaks looks identical, from inside its own measurements, to one that doesn't, leakage can only be surfaced by explicitly regressing one axis's signal on another axis's independently measured features --- turning ``the batteries are orthogonal'' from a silent assumption into a reported, falsifiable number.

\textbf{Method.} Each item's capability-battery delta is regressed on that item's format-feature deltas, and each item's policy-battery delta is regressed on lexical and surface cues, to estimate what fraction of the nominally capability or policy signal is actually explained by surface-form covariates. This explained fraction is reported as an explicit leakage share attached to the pair, rather than being silently absorbed into the capability or policy axis share.

\subsection{Noise-Floor Variance Model}
\textbf{Challenge.} A version-to-version behavioral difference is only evidence of drift if it exceeds what the same version alone would produce across repeated queries --- otherwise ordinary sampling variability, or a silent mid-study infrastructure change, can masquerade as a genuine effect.

\textbf{Insight.} Treating a raw observed delta as if it were already the noise floor is exactly what would let stochastic sampling variability be misread as drift; the floor has to be estimated directly from repeated within-version samples, and its own stability over the study window has to be checked, not assumed.

\textbf{Method.} Each probe item is sampled five times per version at fixed decoding settings to estimate the within-version run-to-run variance for that version and item. A mixed-effects model with axis as a fixed effect and item as a random effect is fit over the resulting per-item deltas, so an axis's estimated contribution to observed drift counts only the portion in excess of the pooled within-version variance floor. Both versions of a pair are queried in interleaved time windows, and the noise-floor estimate from the early half of the study window is compared against the late half to detect any silent infrastructure drift that would otherwise contaminate the floor itself.

\subsection{Cross-Vendor Structural Agreement Test}
\textbf{Challenge.} Whether ``capability,'' ``policy,'' and ``format'' are comparable constructs across different vendors cannot be verified from vendor-internal documentation, because none is available --- the same labeled change might be implemented through entirely different underlying mechanisms from one provider to the next.

\textbf{Insight.} Commensurability across vendors can be tested behaviorally even though it cannot be verified structurally: if an identical instrument, applied identically to every vendor's pairs, produces a rank ordering of axis dominance that agrees across vendors more than a random permutation would, that is evidence the three axes function as a shared construct rather than three vendor-specific artifacts of how the instrument happened to be built.

\textbf{Method.} Each version pair's decomposition is summarized as a point on the three-way simplex of axis shares. Cross-vendor structure is tested two ways: a rank-agreement statistic (Kendall's $W$) computed with pairs as the ranking units --- each of $m \geq 6$ pairs contributes one ranking of the three axes, giving an attainable permutation null ($\chi^2 \approx m(n{-}1)W \approx 7.2$ at $W = 0.60$ with $m = 6$, $n = 3$, $p \approx 0.027$), whereas a vendor-level test with only three vendors as ranking units cannot reach $p < 0.05$ except at perfect agreement ($P(W{=}1) = 1/36 \approx 0.028$) --- and a clustering analysis of the simplex points checking, descriptively, whether pairs group by vendor identity or mix across vendors beyond what the permutation null would predict.

\section{Experiments}

\subsection{Experimental Setup}

\textbf{Datasets.} GSM8K, MMLU-Pro, and HumanEval anchor the capability battery with objectively verifiable, execution- or exact-match-scored items. LiveCodeBench is included specifically for its rolling-release design, which prevents a capability-battery item from having differential training-data exposure across the two versions of a pair --- a confound that would otherwise make a measured capability-axis change indistinguishable from a contamination artifact. XSTest, OR-Bench-Hard-1K, and SORRY-Bench stratify the policy battery into safe-but-seemingly-unsafe, pseudo-harmful, and genuinely harmful prompts, so a refusal-rate change can be read against a fixed ground-truth harmfulness reference rather than confounded with correctness. IFEval anchors the format battery's instructed-format-compliance measurement, alongside a separately pre-verified pool of prompts on which both versions are already near-ceiling correct. Because 2026 frontier pairs sit at or near ceiling on much of GSM8K, HumanEval, and MMLU-Pro --- and a ceiling item carries no capability-drift signal, which would mechanically deflate the capability axis share and inflate the relative format and policy shares, a bias \emph{toward} this paper's own preferred conclusion --- the capability battery applies an off-ceiling pre-verification step that mirrors the format battery's near-ceiling screening in the opposite direction: capability items are pre-screened per pair to a pool on which at least one version is measurably below ceiling, so battery composition is endogenous to the pair rather than to a fixed benchmark's original difficulty. Item-level reliable-change analysis reporting that a majority of MMLU-Pro items are floor- or ceiling-saturated for modern model pairs \citep{BeyondTheMean7405} is the empirical basis for this screening step, and the direction of the residual bias is acknowledged explicitly wherever the format share is headlined.

\textbf{Baselines} (sanity-check / misattribution-demonstration comparators, not competitors to beat --- this is a measurement study, and no research question here is framed as ``our procedure beats baseline X''). The blended-delta audit reproduces the design of the canonical ChatGPT behavior-over-time study \citep{ChatGPTBehavior2307} --- one aggregate per-task delta, no format-robust scoring, no noise floor, no decomposition --- run on the format-only constructed control pairs to check whether standard practice actually misattributes pure format drift as a capability or policy change. The null-pair placebo queries the identical snapshot twice in disjoint interleaved time windows and runs it through the full pipeline; because it is, by construction, the same model against itself, any axis ``detection'' on it is a pipeline artifact rather than evidence of a real phenomenon. A two-way ablation drops the format axis entirely (capability and policy only) and is run on the format-only control pairs to check that the third axis is load-bearing rather than decorative.

\textbf{Metrics.} Axis share $\eta^2_a(P)$ with bootstrap confidence intervals; leakage share $\lambda(P)$; within-version noise floor $\sigma^2_w(P)$ and its early-half-vs-late-half stability; dominant-axis recovery accuracy $\mathrm{Acc}$ on control pairs; Kendall's $W$ with a permutation-null $p$-value for cross-vendor rank agreement.

\textbf{Research questions:}
\begin{enumerate}[leftmargin=1.6em]
\item RQ1 --- Are the probe batteries orthogonal by construction, as judged on format-only control pairs?
\item RQ2 --- Does the full pipeline recover the true dominant axis, blind, on constructed control pairs of known cause?
\item RQ3 --- How large is the residual cross-axis leakage once the batteries and calibration are in place?
\item RQ4 --- Is the estimated within-version noise floor stable enough, across the study window, to serve as a reference?
\item RQ5 --- In real version pairs, do multiple axes exceed the noise floor, and does format carry a non-negligible share?
\item RQ6 --- Does the rank ordering of axis dominance agree across vendors more than chance would predict?
\end{enumerate}

\subsubsection{RQ1: Are the probe batteries orthogonal by construction, as judged on format-only control pairs?}
\textbf{Setup.} On format-only constructed control pairs (identical weights, only the system template altered), measure the capability battery's per-item correctness delta and the policy battery's refusal-rate delta, and compare each against its own within-version noise floor via a permutation test.

\textbf{Expected outcome.} Neither the capability battery's correctness delta nor the policy battery's refusal-rate delta is significantly different from its own within-version noise floor on these format-only pairs.

\textbf{Falsification condition.} Refuted if either the capability or the policy battery shows a delta significantly exceeding its own noise floor on a pair where, by construction, only the format axis changed --- that would mean the battery is not orthogonal to format by design, and the leakage-regression estimate (RQ3) rather than an assumption of orthogonality must be relied on for any real-pair interpretation.

\textbf{Hypothesis (plain language).} The capability and policy batteries do not react to a pure formatting change, as they are designed to.

\textbf{Metrics.} Per-battery permutation-test $p$-value against the within-version noise floor, on format-only control pairs.

\textbf{Result.} \ph

\subsubsection{RQ2: Does the full pipeline recover the true dominant axis, blind, on constructed control pairs of known cause?}
\textbf{Setup.} Across at least nine single-axis constructed control pairs (three per axis) plus at least three mixed-axis control pairs with a known dominant axis (two interventions applied at a known relative dose), run the complete pipeline with pair identities withheld from the analysis code, and compare the predicted dominant axis against the true one, alongside the mean estimated leakage share across the same pairs. The scale-ladder transfer check (Methodology) repeats the format-only control and the null-pair placebo at two or more open-weight scales.

\textbf{Expected outcome.} The pipeline recovers the true dominant axis in at least eight of nine single-axis control pairs and in at least two of three mixed-axis pairs, with mean leakage share below 0.15, and with recovery accuracy and leakage share stable across the scale ladder.

\textbf{Falsification condition.} Refuted if recovery accuracy is not significantly above the one-in-three chance level (exact binomial test, e.g.\ five or fewer of nine correct) or if mean leakage share meets or exceeds the largest reported axis share --- either would mean the instrument is not reliable enough to license any claim about real, uncalibratable pairs, and the paper's central claims (RQ5, RQ6) would have to be reported as unsupported by the calibration gate rather than interpreted at face value. Refuted \emph{for the mixed regime specifically} --- independently of the single-axis result --- if the pipeline recovers pure single-axis causes but fails on the mixed-axis pairs, in which case real-pair decompositions are reported as calibrated only for the pure regime, a substantially weaker license that must be stated alongside every downstream claim; and refuted \emph{for calibration transfer} if the scale ladder shows a strong scale trend in recovery accuracy or leakage share, in which case no calibration number is carried across to commercial-scale pairs without an explicit transfer caveat. Scoring is subject to the pre-registered power rule (Assumption A8, gate PG-7): if the pilot-recomputed power for detecting leakage below 0.15 is under 0.9 for a given axis's three control pairs, that axis's contribution to recovery accuracy is reported as INCONCLUSIVE rather than scored pass/fail, and the eight-of-nine bar is evaluated only over the powered axes with the reduced denominator stated explicitly.

\textbf{Hypothesis (plain language).} The decomposition procedure is calibrated: when the true cause is known, the procedure finds it.

\textbf{Metrics.} Recovery accuracy (exact binomial test against 1/3 chance); mean leakage share with bootstrap confidence interval.

\textbf{Result.} \ph

\subsubsection{RQ3: How large is the residual cross-axis leakage once the batteries and calibration are in place?}
\textbf{Setup.} Regress each control pair's capability-battery and policy-battery deltas on the other axes' independently measured features (format features for capability; lexical/surface cues for policy), and report the explained-variance fraction as the leakage share for that pair.

\textbf{Expected outcome.} The estimated leakage share averages below 0.15 across the constructed control pairs --- the same threshold the calibration gate (RQ2) requires for the decomposition to be considered valid.

\textbf{Falsification condition.} Refuted if the average leakage share meets or exceeds 0.15, or if leakage on any single axis is large enough that dropping that axis's estimate reverses the dominant-axis call on a control pair whose true cause is already known --- either would mean the reported axis shares cannot be interpreted independently of one another.

\textbf{Hypothesis (plain language).} Cross-axis contamination between the batteries is real but small enough not to undermine the reported shares.

\textbf{Metrics.} Mean and per-pair leakage share, with bootstrap confidence intervals.

\textbf{Result.} \ph

\subsubsection{RQ4: Is the estimated within-version noise floor stable enough, across the study window, to serve as a reference?}
\textbf{Setup.} For every tested version pair, compare the within-version noise-floor estimate computed from the early half of the study window against the estimate from the late half.

\textbf{Expected outcome.} The early-half and late-half noise-floor estimates agree within 20\% for every tested pair.

\textbf{Falsification condition.} Refuted if the early-half and late-half estimates diverge by more than 20\% for a pair --- that would indicate a silent mid-study infrastructure change contaminating the noise floor, and any axis-share estimate for that specific pair would have to be reported as unreliable rather than compared against a moving reference.

\textbf{Hypothesis (plain language).} The noise floor is a stable enough reference to compare real drift against, for the duration of the study.

\textbf{Metrics.} Ratio of early-half to late-half noise-floor estimate, per pair and per battery.

\textbf{Result.} \ph

\subsubsection{RQ5: In real version pairs, do multiple axes exceed the noise floor, and does format carry a non-negligible share?}
\textbf{Setup.} Gated by RQ2 passing. Across at least six real version pairs, permutation-test each axis's contribution against its pair's noise floor (Holm--Bonferroni corrected across the pair-by-axis family), and report the format axis's share of explained drift for each pair.

\textbf{Expected outcome.} In at least four of six real version pairs, at least two of the three axes' contributions significantly exceed the noise floor, and the format axis's share is at least 10\% of explained drift in at least three of six pairs.

\textbf{Falsification condition.} Refuted toward ``format is negligible'' if the format axis's share is below 10\% of explained drift in every tested pair once scoring is made format-robust --- in which case the finding becomes that blended audits' apparent format sensitivity is a scoring artifact fixable with format-robust extractors, not evidence of a genuine third behavioral axis, and the paper's Claim 3 is reversed rather than confirmed. Refuted toward ``the decomposition doesn't bind'' if leakage plus noise absorb the majority of the observed difference in most pairs, leaving no axis clearly above its floor. Both the expected outcome and either refutation are scored under the pre-registered power rule (Assumption A8, gate PG-7): any of the eighteen pair-by-axis cells whose pilot-recomputed power is below 0.9 is excluded from both the numerator and the denominator of the ``at least four of six'' and ``at least three of six'' counts and reported separately as INCONCLUSIVE --- never as a pass, a fail, or a silent drop.

\textbf{Hypothesis (plain language).} Version updates typically move more than one axis at once, and format drift --- invisible to a blended score --- is not a negligible contributor to what gets reported as behavioral change.

\textbf{Metrics.} Per-pair, per-axis permutation-test $p$-value against the noise floor; per-pair format-axis share with bootstrap confidence interval.

\textbf{Result.} \ph

\subsubsection{RQ6: Does the rank ordering of axis dominance agree across vendors more than chance would predict?}
\textbf{Setup.} Gated by RQ2 passing and RQ5 binding. Across at least six real version pairs from at least three vendors, compute Kendall's $W$ over the per-pair rankings of the three axis shares --- each pair contributes one ranking as a ranking unit --- tested against a 10{,}000-draw permutation null; separately, check descriptively whether the pairs cluster by vendor identity on the axis-share simplex beyond what the same permutation null would produce. The ranking unit is the pair, not the vendor, because the significance threshold below is otherwise unattainable at the minimum design: with three vendors as ranking units over three axes, the exact permutation null gives $P(W{=}1) = 1/36 \approx 0.028$ and every $W < 1$ has $p > 0.05$, whereas with $m = 6$ pairs, $\chi^2 \approx m(n{-}1)W = 7.2$ at $W = 0.60$, giving $p \approx 0.027$ --- the minimum attainable $p$ at this design is stated alongside the result.

\textbf{Expected outcome.} Kendall's $W$ over per-pair rankings is at least 0.60 with permutation $p < 0.05$, and pairs do not cluster tightly by vendor identity beyond the permuted-null expectation (descriptive clustering check).

\textbf{Falsification condition.} Refuted if Kendall's $W$ is statistically indistinguishable from the permutation null ($p \geq 0.05$), or if pairs cluster tightly by vendor identity beyond the null --- either would indicate the axis-dominance structure is vendor-idiosyncratic rather than shared, meaning longitudinal-audit practice cannot assume one vendor's recipe transfers to another and must budget for a bespoke instrument per vendor. This is a fully publishable answer either way, since the underlying question is two-sided by construction. Dominance ranks entering Kendall's $W$ inherit the same pre-registered power rule (Assumption A8, gate PG-7): a pair whose dominant-axis call rests on a cell reported INCONCLUSIVE is itself reported as INCONCLUSIVE for this test --- excluded from both the agreement statistic and its refutation, never silently dropped.

\textbf{Hypothesis (plain language).} The same three-axis structure that governs one vendor's version updates also governs another's, rather than each vendor having its own idiosyncratic pattern of what a version bump typically changes.

\textbf{Metrics.} Kendall's $W$ (pairs as ranking units) with permutation-null $p$-value and the minimum attainable $p$ at the realized $m$; silhouette score of vendor-labeled clustering against the permuted-label null (descriptive).

\textbf{Result.} \ph

\vspace{0.5em}
\noindent\textbf{Further evaluation.} With more time, a natural extension is to widen the probed behavioral subspace beyond math/knowledge/code capability, harmfulness-stratified refusal, and surface-style format --- to multi-turn, multilingual, or agentic/tool-use behavior, where the three-axis taxonomy itself may need revision rather than simple re-application. This is explicitly out of scope here because it changes what counts as a probe item at a more basic level than extending the current batteries would require.

\section{Related Work}

The motivating predecessor of this paper is the ChatGPT behavior-over-time study \citep{ChatGPTBehavior2307}, and it is worth stating the delta plainly before surveying the wider landscape: that study already documents the exact phenomenon this paper is built to instrument --- a single blended per-task delta between snapshots, whose own reported code-generation drop was later attributable in part to a formatting change scored as a capability loss --- but it offers no mechanism to attribute a detected delta to capability, policy, or format, audits only one vendor, and cannot speak to whether its finding generalizes across providers. This paper's contribution is precisely a calibrated instrument for the attribution that study's own design cannot perform. The nearest \emph{contemporary} prior work, however, is no longer that 2023 audit but the 2026 update-governance and change-detection cluster surveyed next, and the positioning below concedes each overlap explicitly rather than claiming an empty field.

\textbf{Detecting and governing model updates without attribution.} A 2026 line of work is this paper's nearest contemporary neighborhood. Deployer-side update governance \citep{TestBeforeYouDeploy7789} tests an updated model by risk categories --- functionality, formatting, and safety among them, the same trichotomy measured here --- but as per-category pass/fail acceptance gates, and it explicitly lists detecting and explaining model drift under limited provider transparency as an open challenge rather than a solved problem; this paper is an instrument for exactly that attribution step. Black-box change detection \citep{TokenEfficientChangeDet1083} flags that a served model differs from a reference using minimal queries, and item-level reliable-change analysis \citep{BeyondTheMean7405} identifies which benchmark items genuinely changed against a within-model sampling-variability floor --- the latter anticipates the within-version noise floor used here, so the surviving delta on that component is the noise floor's role inside a three-way attribution, together with the separately reported cross-axis leakage share, rather than binary change detection. Refusal--capability separation via abliteration \citep{WillingButUnable5396} cleanly disentangles two of this paper's three axes, but through a white-box weight-space intervention unavailable for API-served pairs and with no version-pair framing. Anytime-valid attribution of evaluation-pipeline drift \citep{WhoDrifted5474} apportions a metric shift between the serving system and the judge with a calibrated sequential test --- the closest methodological relative of this paper's calibration-first stance --- but it attributes across a different pair of causes, relies on a human-labeled anchor set, and has no capability/policy/format decomposition. Across this cluster, updates are detected, gated, or two-way attributed; none performs a black-box three-way attribution of a version pair's drift, and none validates its attribution blind on constructed pairs of known cause before interpreting real ones.

\textbf{Longitudinal LLM behavior auditing.} The ChatGPT behavior-over-time study \citep{ChatGPTBehavior2307} is the canonical demonstration that deployed model behavior changes measurably across shipped snapshots, evaluating a fixed task suite spanning math, sensitive questions, opinion surveys, code generation, and reasoning. Work on adapter robustness across evolving base-model checkpoints \citep{CARLoRA5Gim} documents a downstream consequence of the same phenomenon: adapters trained against one snapshot silently degrade as the underlying model updates, forcing full retraining. Both establish that version drift is real and consequential, but both treat it as an undifferentiated before/after or before/after-training perturbation rather than something to be apportioned across competing causes --- the former reports one blended number per task, and the latter trains against drift as a nuisance to be robust to rather than a phenomenon to diagnose.

\textbf{Attributing behavioral change with model-internal access.} Stage-wise alignment-impact analysis \citep{MoreRLHFMoreTrustFpiC} attributes behavioral differences between training checkpoints to specific stages of a training pipeline using influence-function-based data attribution, but requires the checkpoints and training data of a controlled, white-box pipeline that black-box commercial version pairs never expose, and its measured dimensions are all trustworthiness facets with no capability or format axis. Diagnosing shallow alignment via per-token KL divergence between an aligned model and its base \citep{FewTokensDeep6Mxh} isolates exactly where in a response safety effects concentrate, but requires logit access to both models being compared. Learning to translate a weight difference between two models into a natural-language description of the behavioral change it induces \citep{DiffInterpretation6As4} likewise requires full weight access to both models. All three demonstrate that fine-grained attribution is possible when internals are available, and none applies to the exact pairs --- commercial API snapshots --- where auditors most need an answer.

\textbf{Separating policy from capability and format at a single time point.} A rank-correlation measure separating a model's refusal probability from its error probability on factual questions \citep{RefusalIndex9gJB} demonstrates, within a single model, that refusal behavior (policy) and knowledge limitation (capability) are distinct, separately measurable constructs --- but it is a one-time-point calibration measure with no version-pair or drift framing, and no format axis. Work on reducing over-refusal \citep{ContrastiveOverRefusalGXcN} shows that safety fine-tuning couples truly harmful and merely superficially risky prompts in a model's internal representation space, providing empirical grounding for treating benign-but-superficially-risky prompts as a discriminative stratum for a policy probe --- and its safety fine-tuning recipe is the one this paper's own policy control-pair construction (Methodology) adapts --- but its diagnosis requires representation-space access and its goal is fixing one model, not measuring drift between two deployed versions. Judge-based evaluation frameworks \citep{WildBenchMKEH} document that automated scoring is systematically biased by response length and style, and correct for this with length-penalized metrics inside a single-time-point leaderboard --- evidence that surface-form confounds are a real, previously documented measurement problem, but treated there as a nuisance to control away rather than a first-class, separately reported behavioral axis (and, notably, a length-specific correction rather than a general theory of format-robust scoring across the markdown/list/hedging axes this paper's own format battery targets).

\section{Conclusion}

This work reframes a widespread but under-examined assumption in longitudinal LLM auditing --- that a blended before/after score is an adequate description of what changed between two model versions --- as a measurement problem with a precise identifiability structure, rather than treating the blended number as self-evidently interpretable. The central methodological wager is that capability, policy, and format drift are only separable once cross-axis leakage and within-version noise are estimated and reported as first-class quantities in their own right, and once the entire procedure is calibrated on pairs of known cause before any claim is made about pairs where the cause is unverifiable; a weaker design risks reporting a number that looks precise while silently conflating exactly the sources this paper is built to keep apart. Whichever result the calibration gate and the real-pair decomposition ultimately produce, the finding is decision-relevant either way: a validated, non-negligible format share would mean current audit practice is systematically over-crediting or over-blaming capability and policy for what is sometimes just a formatting change, while a negative calibration result would itself be an important finding about the limits of what black-box behavioral measurement can distinguish. Three limitations bound this claim from the outset and are natural next steps rather than afterthoughts: attribution for commercial, black-box version pairs remains a calibrated behavioral inference that is never verified against vendor internals, since no method here resolves that opacity, only bounds the uncertainty it creates; the cross-vendor structural claim rests on a small number of pairs and vendors, so its result is reported as graded statistical evidence with a stated permutation-null comparison, not a binary claim about the industry as a whole; and the calibration gate itself is bounded by its own construction --- its control pairs are built under the very taxonomy they validate and on open-weight models smaller than the commercial stacks the instrument targets, so the mixed-axis controls and the scale-ladder transfer check bound, but cannot eliminate, the residual risk that calibration success on constructed pairs overstates reliability on real ones.

\vspace{1em}
\hrule
\vspace{1em}

\section*{Pre-handoff checklist}
\begin{itemize}[leftmargin=1.4em]
\item \textbf{Organization/navigation} --- narrative flows problem $\to$ method $\to$ experiments $\to$ related work without forward references to undefined terms.
\item \textbf{Terminology discipline} --- ``axis share,'' ``leakage share,'' ``noise floor,'' ``dominant axis,'' ``recovery accuracy'' each defined at first use in Problem Definition, used consistently thereafter; no internal stage IDs (\texttt{sp\_*}, \texttt{m\_*}, \texttt{rq\_*}) leaked into prose --- checked directly against the discussion pipeline log.
\item \textbf{Notation with worked example} --- the formal setup section defines $(P, v_{\mathrm{old}}, v_{\mathrm{new}}, a, \eta^2_a, \lambda, \sigma^2_w, \mathrm{dom}(P), \mathrm{Acc}, W)$ and the Problem Definition/Methodology sections consistently reuse this notation; no separate worked numeric example is included since no numbers exist yet (would require fabricating results).
\item \textbf{RQ framing matches outcomes} --- all 6 RQs are graded/quantitative (proportions, agreement rates, statistical thresholds), each answered in its own subsection with an explicit falsifiable expected outcome and a distinct falsification condition, not a binary yes/no.
\item \textbf{No typesetting artifacts} --- only explicit placeholder markers in Results, per the drafting convention this is a pre-registered spec, not a report.
\item \textbf{Artifact plan} --- the three probe batteries, the control-pair construction recipes, and the cached raw outputs from every audited version pair are the natural release artifacts; not yet built, so not claimed as available.
\end{itemize}

\section*{Hand-off}
This is a draft, not a finished paper --- experiments have not run. Natural next steps: prose polish and venue-specific LaTeX formatting once a target venue is chosen, or revising an earlier stage if piloting surfaces a problem (e.g., if recovery accuracy on control pairs is far below the pre-registered gate).

\bibliographystyle{plainnat}
\bibliography{refs}

\end{document}
